% Paragrafo con titolo in grassetto e corpo rientrato.
\newenvironment{paragrafo}[1]{%
    \par
    \noindent\textbf{#1}\par
    \begin{adjustwidth}{2em}{0pt}%
    \setlength{\parindent}{0pt}%
    \noindent\ignorespaces
}{%
    \end{adjustwidth}%
    \par
}

% Ambienti semantici con la stessa resa grafica di paragrafo.
\newenvironment{teorema}
    {\begin{paragrafo}{Teorema.}}
    {\end{paragrafo}}

\newenvironment{dimostrazione}
    {\begin{paragrafo}{Dimostrazione}}
    {\end{paragrafo}}

% Stile condiviso dai diagrammi di memoria.
\tikzset{
    cella memoria/.style={
        draw,
        minimum width=0.8cm,
        minimum height=0.55cm,
        inner sep=0pt
    }
}

\newcommand{\memoria}[1]{%
\begin{tikzpicture}
\foreach \val [count=\i from 0] in {#1}{
    \node[cella memoria] (c\i) at (\i*1.05,0) {\val};
    \node[above=2pt of c\i] {\scriptsize \i};
}
\end{tikzpicture}%
}

\newcommand{\memoriaunita}[1]{%
\begin{tikzpicture}
\foreach \val [count=\i from 0] in {#1}{
    \node[cella memoria, outer sep=0pt, anchor=west]
        (c\i) at ({0.8*\i},0) {\val};
    \node[above=2pt of c\i] {\scriptsize \i};
}
\end{tikzpicture}%
}
